\documentclass[11pt,a4paper]{article}
\usepackage[utf8]{inputenc}
\usepackage[T1]{fontenc}
\usepackage{amsmath,amssymb,amsthm}
\usepackage{listings}
\usepackage{geometry}
\usepackage{hyperref}
\geometry{margin=2.5cm}

\newtheorem{lemma}{Lemma}
\newtheorem{definition}{Definition}

\lstset{
  language=C++,
  basicstyle=\ttfamily\small,
  keywordstyle=\bfseries,
  breaklines=true,
  numbers=left,
  numberstyle=\tiny,
  frame=single,
  tabsize=2
}

\title{IOI 2019 -- Rectangles}
\author{Solution}
\date{}

\begin{document}
\maketitle

\section{Problem Summary}
Given an $n \times m$ grid of distinct integers, count rectangles $(r_1, r_2, c_1, c_2)$ with $r_1 < r_2$, $c_1 < c_2$ such that every boundary cell is strictly greater than every interior cell.

\section{Solution}

\subsection{Two Independent Conditions}
A rectangle $(r_1, r_2, c_1, c_2)$ is valid iff:
\begin{enumerate}
  \item \textbf{Row condition}: For each interior column $c \in (c_1, c_2)$, $\max(a[r_1+1..r_2-1][c]) < \min(a[r_1][c], a[r_2][c])$.
  \item \textbf{Column condition}: For each interior row $r \in (r_1, r_2)$, $\max(a[r][c_1+1..c_2-1]) < \min(a[r][c_1], a[r][c_2])$.
\end{enumerate}

\subsection{Sparse Tables for Range Max}
Precompute sparse tables for each column (range max over rows) and each row (range max over columns), enabling $O(1)$ queries.

\subsection{Enumeration}
For each pair of rows $(r_1, r_2)$, scan interior columns to find maximal runs satisfying the row condition. For each candidate column pair $(c_1, c_2)$ from such a run, verify the column condition.

\section{Implementation}

\begin{lstlisting}
#include <bits/stdc++.h>
using namespace std;

long long count_rectangles(vector<vector<int>> a) {
    int n = a.size(), m = a[0].size();
    if (n < 3 || m < 3) return 0;

    int LOG = 1;
    while ((1 << LOG) < max(n, m)) LOG++;

    // Sparse table: column c, range max over rows
    vector<vector<vector<int>>> colMax(m, vector<vector<int>>(LOG, vector<int>(n)));
    for (int c = 0; c < m; c++) {
        for (int r = 0; r < n; r++) colMax[c][0][r] = a[r][c];
        for (int k = 1; k < LOG; k++)
            for (int r = 0; r + (1 << k) <= n; r++)
                colMax[c][k][r] = max(colMax[c][k-1][r],
                                       colMax[c][k-1][r + (1 << (k-1))]);
    }
    auto qcol = [&](int c, int lo, int hi) -> int {
        if (lo > hi) return INT_MIN;
        int k = __lg(hi - lo + 1);
        return max(colMax[c][k][lo], colMax[c][k][hi - (1 << k) + 1]);
    };

    // Sparse table: row r, range max over columns
    vector<vector<vector<int>>> rowMax(n, vector<vector<int>>(LOG, vector<int>(m)));
    for (int r = 0; r < n; r++) {
        for (int c = 0; c < m; c++) rowMax[r][0][c] = a[r][c];
        for (int k = 1; k < LOG; k++)
            for (int c = 0; c + (1 << k) <= m; c++)
                rowMax[r][k][c] = max(rowMax[r][k-1][c],
                                       rowMax[r][k-1][c + (1 << (k-1))]);
    }
    auto qrow = [&](int r, int lo, int hi) -> int {
        if (lo > hi) return INT_MIN;
        int k = __lg(hi - lo + 1);
        return max(rowMax[r][k][lo], rowMax[r][k][hi - (1 << k) + 1]);
    };

    long long answer = 0;
    for (int r1 = 0; r1 < n - 2; r1++) {
        for (int r2 = r1 + 2; r2 < n; r2++) {
            // Find maximal runs of valid interior columns
            int runStart = -1;
            for (int c = 1; c < m - 1; c++) {
                int imax = qcol(c, r1 + 1, r2 - 1);
                bool ok = (imax < a[r1][c] && imax < a[r2][c]);
                if (!ok) { runStart = -1; continue; }
                if (runStart == -1) runStart = c;
                // Check all column pairs ending at c within this run
                for (int c1 = runStart - 1; c1 < c; c1++) {
                    int c2 = c + 1;
                    // Verify c1 >= 0, c2 <= m-1 (c1 = boundary col, c2 = boundary col)
                    if (c1 < 0 || c2 >= m) continue;
                    // Verify all interior cols c1+1..c2-1 are in run
                    if (c1 + 1 < runStart) continue;
                    // Check column condition for all interior rows
                    bool colOk = true;
                    for (int r = r1 + 1; r <= r2 - 1 && colOk; r++) {
                        int rm = qrow(r, c1 + 1, c2 - 1);
                        if (rm >= a[r][c1] || rm >= a[r][c2]) colOk = false;
                    }
                    if (colOk) answer++;
                }
            }
        }
    }
    return answer;
}

int main() {
    int n, m;
    scanf("%d %d", &n, &m);
    vector<vector<int>> a(n, vector<int>(m));
    for (int i = 0; i < n; i++)
        for (int j = 0; j < m; j++)
            scanf("%d", &a[i][j]);
    printf("%lld\n", count_rectangles(a));
    return 0;
}
\end{lstlisting}

\section{Complexity Analysis}
\begin{itemize}
  \item \textbf{Preprocessing}: $O(nm \log(\max(n,m)))$ for sparse tables.
  \item \textbf{Enumeration}: $O(n^2 m^2)$ worst case. The full-score solution uses monotone stacks to enumerate valid pairs efficiently, achieving $O(n^2 m + nm^2)$.
  \item \textbf{Space}: $O(nm \log(\max(n,m)))$.
\end{itemize}

\textbf{Note.} For $n, m \le 2500$, the intended solution uses monotone stacks on each row and column to enumerate valid $(r_1, r_2)$ and $(c_1, c_2)$ pairs, combined with 2D prefix sums for efficient intersection counting.

\end{document}
